\section{Acknowledgement of generative-AI tools}
\label{sec:ai}

This section is required by the course rubric and is a compressed version of
\texttt{AI\_USAGE.md} at the repository root, which carries the full per-tool table
and policy statement.

\subsection*{Team policy}

The team used generative-AI tools as collaborators, not as authors. Three
durable rules governed every contribution:

\begin{enumerate}
  \item Every line of code that lands in \texttt{main} is reviewed by a human. AI
        output that no team member can explain is not merged.
  \item No AI tool is given access to credentials, API keys, or production data.
        The \texttt{.env} file is gitignored; only \texttt{.env.example} (no secrets)
        ships in the repository.
  \item The LLM features inside the product itself are opt-in:
        \texttt{backend/config.py} defaults to \texttt{LLM\_ENABLED=false}, the
        test suite uses the deterministic \texttt{NullProvider}, and activating
        the live \texttt{GeminiProvider} requires three environment variables.
\end{enumerate}

We take collective responsibility for the whole repository, including text and
code that was originally drafted by an AI tool and then reviewed and accepted.

\subsection*{Tools used}

\textbf{Claude Code} (Anthropic, CLI agent) was the primary coding assistant for
multi-file refactors, agent scaffolding, documentation drafts, and the bulk of the
LaTeX source for this submission; it left footprints in the repository via the
\texttt{.github/workflows/claude.yml} (issue/PR \texttt{@claude} mention handler)
and \texttt{.github/workflows/claude-code-review.yml} (per-PR \texttt{/code-review})
workflows. \textbf{Claude.ai} (Anthropic, web) was used by team members for design
discussions and prose drafting. \textbf{ChatGPT} (OpenAI) was used opportunistically
for isolated syntax checks and LaTeX-table formatting. \textbf{GitHub Copilot}
(GitHub/OpenAI) provided in-editor autocomplete for routine code; one merge-resolution
commit (\texttt{b9fc033}) was authored by the \texttt{copilot-swe-agent[bot]} account
as a one-off, not a sustained authoring pattern. \textbf{Cursor} (Anysphere) was used
by one team member as their primary editor for some of the frontend work.
\textbf{Google Gemini} (web) was used to spot-check factual claims about NYSE hours
and \texttt{ccxt} pagination semantics; \textbf{Google Gemini API} is also wired into
the product itself as the optional live \texttt{LLMProvider}, but that is a product
feature, not a development tool.

The team did \emph{not} use Codex, Replit Ghostwriter, Sourcegraph Cody, Tabnine,
Codeium, or Amazon CodeWhisperer.

\subsection*{Where AI was not used}

The strategy formulas in \texttt{backend/strategies/*.py} were transcribed from the
academic and practitioner sources cited in \texttt{CITATIONS.md} and
Section~\ref{sec:methodology}; the team independently re-derived each one before
implementation. The look-ahead-bias enforcement in
\texttt{backend/backtest/engine.py} and its oracle-strategy test were designed and
written by the team. The database schema, the API contract, and the four-week
project timeline were team decisions, not AI-generated.

\subsection*{Reproducibility caveat}

AI sessions are not natively reproducible across runs or prompt phrasings. We
mitigate this via diff-level human review, the test suite's coverage of the
critical correctness invariants, and prompt summaries in the project diary
(Section~\ref{sec:diary}) for the most impactful changes.
